% GENERATED FILE. Edit canonical lessons, then run npm run generate:books.
% canonical-source-hash: fnv1a64:a03020a1adfb30b5
% canonical-lessons: ZH-C228-touteng, ZH-W228-tou-head, ZH-W228-teng-ache, ZH-C228-chiyao, ZH-C228-dazhen, ZH-C228-yanzhong, ZH-C228-diantou

\chapter{A Headache}
\label{ch:zh-a-headache}
\hlchaptermodality{228}

\begin{chapteropening}
\textbf{By the end of this chapter:} \emph{I can say my head aches, take medicine, have an injection, say whether it is serious, and nod.}

Complete the last lesson of chapter 228: I can say my head aches, take medicine, have an injection, say whether it is serious, and nod.
\end{chapteropening}

\section[\texorpdfstring{\zh{头疼} (tóuténg)}{头疼 (tóuténg)}]{\zh{头疼} (tóuténg) — having a headache; aching in the head}

\label{lesson:ZH-C228-touteng}



\begin{quote}
Before the new one: say the Mandarin for a sandy beach, then the Mandarin for to run into someone.
\end{quote}

\subsection*{You'll want to know: \zh{头疼}}
\textbf{\zh{头疼}} — \emph{tóuténg} — \textquotedblleft{}having a headache; aching in the head\textquotedblright{}.

\begin{quote}
\textbf{\zh{我头疼不想吃饭}} \emph{wǒ tóuténg bù xiǎng chīfàn} — I've got a headache, so I don't feel like eating.
\end{quote}

The new character, in pieces: \zh{头} — five strokes: two dots, a long bar, a long stroke falling to the left, and a dot at the lower right; \zh{疼} — ten strokes: the illness frame of \zh{病}, five strokes, then inside it five more: a short fall to the left, a bar that turns and falls to the left, a long stroke to the right, and two dots. The next lesson puts it in your hand.

\begin{cousinweb}
\textbf{\zh{头疼}} (\emph{tóuténg}) is built from \zh{头} 'head' and \zh{疼} 'sore', both inherited from Middle Chinese, \emph{duw} and \emph{downg} (Guangyun, in Baxter's transcription), which glosses the first 'the head' and the second 'pain'.
\end{cousinweb}

\begin{grammarlens}[title={What you've built}]
One more word for giving reasons.
\end{grammarlens}

\subsection*{Your turn}
\begin{itemize}
\raggedright
  \item \emph{Say it:} \emph{tóuténg}
  \item \emph{Say it:} \emph{tóuténg}, once more
  \item \emph{Say it:} \emph{wǒ tóuténg bù xiǎng chīfàn} — I've got a headache, so I don't feel like eating.
  \item \emph{Recall:} say the Mandarin for a sandy beach, then the Mandarin for to run into someone, then say \emph{tóuténg} again
\end{itemize}

\begin{tcolorbox}[breakable,colback=teal!4,colframe=teal!35!black,title={Before you move on}]
What does \zh{头疼} mean? (\textquotedblleft{}Having a headache; aching in the head\textquotedblright{}.) Say it once more.
\end{tcolorbox}
\section[\texorpdfstring{\zh{头} (tóu)}{头 (tóu)}]{\zh{头} — two dots over a shape like \zh{大}}

\label{lesson:ZH-W228-tou-head}



\begin{quote}
The word first, then the character inside it.

\begin{itemize}
\raggedright
  \item \emph{Recall:} say \emph{tóuténg} — \textbf{R1}, one lesson back, before it decays
\end{itemize}

\begin{quote}
\zh{大} \zh{天}
\end{quote}

\emph{Read it:} each one aloud, with nothing said first — \emph{dà}, \emph{tiān}
\end{quote}

\subsection*{Script — two dots over a shape like \zh{大}}
\hlblockfigure{\detokenize{figures/ZH-W228-tou-head-filmstrip.pdf}}{How \zh{头} is written, stroke by stroke}

\begin{quote}
\zh{头}
\end{quote}

\textbf{Two pieces: two dots, then a shape like \zh{大}.}

\textbf{Two dots,} one under the other, each falling to the right, a little left of the middle.

\textbf{Then three strokes, almost \zh{大}:} a long bar across, then a long stroke from the top falling down to the left. The last stroke is not the long stroke of \zh{大} but a short dot at the lower right.

2 + 3 = five.

\subsection*{Your turn}
\begin{itemize}
\raggedright
  \item \emph{Write it:} \zh{头}, all five, with the model in front of you
  \item \emph{Cover:} write it again from memory, and count the pieces, not the strokes
\end{itemize}

\begin{tcolorbox}[breakable,colback=teal!4,colframe=teal!35!black,title={Before you move on}]
How many strokes in \zh{头}? (\textbf{Five}.) What comes first? (\textbf{Two dots}.) How is the last stroke different from \zh{大}? (\textbf{It is a short dot}.)

Source: \href{https://raw.githubusercontent.com/chanind/hanzi-writer-data/68d10a4b21150cae5e1ebbd223eed289cf32d90c/data/\%E5\%A4\%B4.json}{Hanzi Writer Data \zh{头}, pinned stroke paths and medians}.
\end{tcolorbox}
\section[\texorpdfstring{\zh{疼} (téng)}{疼 (téng)}]{\zh{疼} — the frame of \zh{病}, and a shape of five inside}

\label{lesson:ZH-W228-teng-ache}



\begin{quote}
The word first, then the character inside it.

\begin{itemize}
\raggedright
  \item \emph{Recall:} say \emph{tóuténg} — \textbf{R1}, one lesson back, before it decays
\end{itemize}

\begin{quote}
\zh{病}
\end{quote}

\emph{Read it:} each one aloud, with nothing said first — \emph{bìng}
\end{quote}

\subsection*{Script — the frame of \zh{病}, and a shape of five inside}
\hlblockfigure{\detokenize{figures/ZH-W228-teng-ache-filmstrip.pdf}}{How \zh{疼} is written, stroke by stroke}

\begin{quote}
\zh{疼}
\end{quote}

\textbf{Two pieces: a frame, and what sits inside it.} Count the pieces, not the strokes.

\textbf{The frame, five strokes,} as in \zh{病}: a dot at the top, a bar across, a long stroke falling to the left, then a small dot and a short stroke rising to the right on the far left.

\textbf{Inside, five strokes.} A short stroke falling to the left. A bar that turns and falls to the left. A long stroke falling to the right. Then two dots, one under the other.

5 + 5 = ten.

\subsection*{Your turn}
\begin{itemize}
\raggedright
  \item \emph{Write it:} \zh{疼}, all ten, with the model in front of you
  \item \emph{Cover:} write it again from memory, and count the pieces, not the strokes
\end{itemize}

\begin{tcolorbox}[breakable,colback=teal!4,colframe=teal!35!black,title={Before you move on}]
How many strokes in \zh{疼}? (\textbf{Ten}.) Where has your hand made the frame before? (\textbf{In \zh{病}}.) What comes at the foot? (\textbf{Two dots}.)

Source: \href{https://raw.githubusercontent.com/chanind/hanzi-writer-data/68d10a4b21150cae5e1ebbd223eed289cf32d90c/data/\%E7\%96\%BC.json}{Hanzi Writer Data \zh{疼}, pinned stroke paths and medians}.
\end{tcolorbox}
\section[chīyào]{chīyào — to take medicine}

\label{lesson:ZH-C228-chiyao}



\begin{quote}
Before the new one: say the Mandarin for to run into someone, then the Mandarin for a headache.
\end{quote}

\subsection*{You'll want to know: chīyào}
\emph{chīyào} — \textquotedblleft{}to take medicine\textquotedblright{}.

\begin{quote}
\emph{nǐ chīyào le ma} — Have you taken your medicine?
\end{quote}

Its characters come later; for now it is a word for the ear.

\begin{cousinweb}
\emph{chīyào} is built from \zh{吃} 'to eat' and \emph{yào} 'medicine'. The second is inherited from Middle Chinese \emph{yak} (Guangyun, in Baxter's transcription), explained, after an older book, as a plant that cures illness. The dictionary gives the first, \emph{kj+t} there, only for stammering, so how it came to mean eating is not settled, and chī is not the regular outcome of that form.
\end{cousinweb}

\begin{grammarlens}[title={What you've built}]
One more word for giving reasons.
\end{grammarlens}

\subsection*{Your turn}
\begin{itemize}
\raggedright
  \item \emph{Say it:} \emph{chīyào}
  \item \emph{Say it:} \emph{chīyào}, once more
  \item \emph{Say it:} \emph{nǐ chīyào le ma} — Have you taken your medicine?
  \item \emph{Recall:} say the Mandarin for to run into someone, then the Mandarin for a headache, then say \emph{chīyào} again
\end{itemize}

\begin{tcolorbox}[breakable,colback=teal!4,colframe=teal!35!black,title={Before you move on}]
What does chīyào mean? (\textquotedblleft{}To take medicine\textquotedblright{}.) Say it once more.
\end{tcolorbox}
\section[dǎzhēn]{dǎzhēn — to have an injection; to give one}

\label{lesson:ZH-C228-dazhen}



\begin{quote}
Before the new one: say the Mandarin for a headache, then the Mandarin for to take medicine.
\end{quote}

\subsection*{You'll want to know: dǎzhēn}
\emph{dǎzhēn} — \textquotedblleft{}to have an injection; to give one\textquotedblright{}.

\begin{quote}
\emph{wǒ pà dǎzhēn} — I'm scared of injections.
\end{quote}

Its characters come later; for now it is a word for the ear.

\begin{cousinweb}
\emph{dǎzhēn} is built from \zh{打} 'to strike, to do' and \emph{zhēn} 'needle', both inherited from Middle Chinese, \emph{taengX} and \emph{tsyim} (Guangyun, in Baxter's transcription), which glosses the second 'needle and thread'.
\end{cousinweb}

\begin{grammarlens}[title={What you've built}]
One more word for giving reasons.
\end{grammarlens}

\subsection*{Your turn}
\begin{itemize}
\raggedright
  \item \emph{Say it:} \emph{dǎzhēn}
  \item \emph{Say it:} \emph{dǎzhēn}, once more
  \item \emph{Say it:} \emph{wǒ pà dǎzhēn} — I'm scared of injections.
  \item \emph{Recall:} say the Mandarin for a headache, then the Mandarin for to take medicine, then say \emph{dǎzhēn} again
\end{itemize}

\begin{tcolorbox}[breakable,colback=teal!4,colframe=teal!35!black,title={Before you move on}]
What does dǎzhēn mean? (\textquotedblleft{}To have an injection; to give one\textquotedblright{}.) Say it once more.
\end{tcolorbox}
\section[yánzhòng]{yánzhòng — serious (of illness); serious (of a problem)}

\label{lesson:ZH-C228-yanzhong}



\begin{quote}
Before the new one: say the Mandarin for to take medicine, then the Mandarin for to have an injection.
\end{quote}

\subsection*{You'll want to know: yánzhòng}
\emph{yánzhòng} — \textquotedblleft{}serious (of illness); serious (of a problem)\textquotedblright{}.

\begin{quote}
\emph{yīshēng shuō bù yánzhòng} — The doctor says it isn't serious.
\end{quote}

Its characters come later; for now it is a word for the ear.

\begin{cousinweb}
\emph{yánzhòng} is built from \emph{yán} 'stern' and \zh{重} 'heavy', both inherited from Middle Chinese, \emph{ngjaem} and \emph{drjowngX} (Guangyun, in Baxter's transcription), which glosses the first 'stern' and 'awe-inspiring'.
\end{cousinweb}

\begin{grammarlens}[title={What you've built}]
One more word for giving reasons.
\end{grammarlens}

\subsection*{Your turn}
\begin{itemize}
\raggedright
  \item \emph{Say it:} \emph{yánzhòng}
  \item \emph{Say it:} \emph{yánzhòng}, once more
  \item \emph{Say it:} \emph{yīshēng shuō bù yánzhòng} — The doctor says it isn't serious.
  \item \emph{Recall:} say the Mandarin for to take medicine, then the Mandarin for to have an injection, then say \emph{yánzhòng} again
\end{itemize}

\begin{tcolorbox}[breakable,colback=teal!4,colframe=teal!35!black,title={Before you move on}]
What does yánzhòng mean? (\textquotedblleft{}Serious (of illness); serious (of a problem)\textquotedblright{}.) Say it once more.
\end{tcolorbox}
\section[\texorpdfstring{\zh{点头} (diǎntóu)}{点头 (diǎntóu)}]{\zh{点头} (diǎntóu) — to nod}

\label{lesson:ZH-C228-diantou}



\begin{quote}
Before the new one: say the Mandarin for to have an injection, then the Mandarin for serious (of illness).
\end{quote}

\subsection*{You'll want to know: \zh{点头}}
\textbf{\zh{点头}} — \emph{diǎntóu} — \textquotedblleft{}to nod\textquotedblright{}.

\begin{quote}
\textbf{\zh{我问他头疼不疼他点头了}} \emph{wǒ wèn tā tóu téng bu téng tā diǎntóu le} — I asked if his head hurt, and he nodded.
\end{quote}

\begin{cousinweb}
\textbf{\zh{点头}} (\emph{diǎntóu}) is built from \zh{点} 'to dot' and \zh{头} 'head', both inherited from Middle Chinese, \emph{temX} and \emph{duw} (Guangyun, in Baxter's transcription). The dictionary glosses the first 'dots and strokes'; 'to dip' is not given there.
\end{cousinweb}

\begin{grammarlens}[title={What you've built}]
One more word for giving reasons.
\end{grammarlens}

\subsection*{Your turn}
\begin{itemize}
\raggedright
  \item \emph{Say it:} \emph{diǎntóu}
  \item \emph{Say it:} \emph{diǎntóu}, once more
  \item \emph{Say it:} \emph{wǒ wèn tā tóu téng bu téng tā diǎntóu le} — I asked if his head hurt, and he nodded.
  \item \emph{Recall:} say the Mandarin for to have an injection, then the Mandarin for serious (of illness), then say \emph{diǎntóu} again
\end{itemize}

\begin{tcolorbox}[breakable,colback=teal!4,colframe=teal!35!black,title={Before you move on}]
What does \zh{点头} mean? (\textquotedblleft{}To nod\textquotedblright{}.) Say it once more.
\end{tcolorbox}
